\begin{flushright}
  \fechaclase{27 de agosto de 2026}
\end{flushright}

\paragraph{Teorema: el complemento de Schur}
La matriz
\[
  S=
  \begin{bmatrix}
    S_{11} & S_{12}\\
    S_{21} & S_{22}
  \end{bmatrix}
\]
es definida positiva si
\[
  S_{11}>0
\]
\[
  S_{22}>0
\]
\[
  S_{11}-S_{12}S_{22}^{-1}S_{12}^{T}>0
\]
\[
  S_{22}-S_{12}^{T}S_{11}^{-1}S_{12}>0
\]

\medskip
\noindent\textit{\textbf{Demostración.}}
Suponga que $S>0$, entonces para cualquier
\[
  X=
  \begin{bmatrix}
    X_1\\
    X_2
  \end{bmatrix},
  \qquad
  X_1\in\mathbb{R}^{p},
  \qquad
  X_2\in\mathbb{R}^{q}
\]
distintos de cero tenemos
\[
  0<X^{T}SX
  =
  \begin{bmatrix}
    X_1^{T} & X_2^{T}
  \end{bmatrix}
  \begin{bmatrix}
    S_{11} & S_{12}\\
    S_{21} & S_{22}
  \end{bmatrix}
  \begin{bmatrix}
    X_1\\
    X_2
  \end{bmatrix}
\]
\[
  0<X_1^{T}S_{11}X_1
  +X_1^{T}S_{12}X_2
  +X_2^{T}S_{12}X_1
  +X_2^{T}S_{22}X_2
\]
Si $X_2=0\Rightarrow X_2^{T}=0\Rightarrow
X_1^{T}S_{11}X_1>0$ entonces $S_{11}>0$.

De manera análoga

Si $X_1=0\Rightarrow X_2^{T}=0\Rightarrow
X_2^{T}S_{22}X_2>0$ entonces $S_{22}>0$.

Eligiendo
\[
  X_2=-S_{22}^{-1}S_{12}^{T}X_1
\]
\begin{align*}
  0<{}&X_1^{T}S_{11}X_1
  +X_1^{T}S_{12}(-S_{22}^{-1}S_{12}^{T}X_1)
  +(-S_{22}^{-1}S_{12}^{T}X_1)^{T}S_{12}^{T}X_1
  +(-S_{22}^{-1}S_{12}^{T}X_1)^{T}
  S_{22}(-S_{22}^{-1}S_{12}^{T}X_1)
\end{align*}
\begin{align*}
  0<{}&X_1^{T}S_{11}X_1
  -X_1^{T}S_{12}S_{22}^{-1}S_{12}^{T}X_1
  -X_1^{T}S_{12}S_{22}^{-1}S_{12}^{T}X_1
  +X_1^{T}S_{12}S_{22}^{-1}S_{22}
  S_{22}^{-1}S_{12}^{T}X_1
\end{align*}
\[
  0<X_1^{T}S_{11}X_1
  -X_1^{T}S_{12}S_{22}^{-1}S_{12}^{T}X_1
  -X_1^{T}S_{12}S_{22}^{-1}S_{12}^{T}X_1
  +X_1^{T}S_{12}S_{22}^{-1}S_{12}^{T}X_1
\]
\[
  0<X_1^{T}S_{11}X_1
  -X_1^{T}S_{12}S_{22}^{-1}S_{12}^{T}X_1
\]
\[
  0<X_1^{T}
  (S_{11}-S_{12}S_{22}^{-1}S_{12}^{T})X_1
\]
De modo que
\[
  S_{11}-S_{12}S_{22}^{-1}S_{12}^{T}>0
\]

Haciendo
\[
  X_1=-S_{11}^{-1}S_{12}X_2
\]
\begin{align*}
  0<{}&(-S_{11}^{-1}S_{12}X_2)^{T}
  S_{11}(-S_{11}^{-1}S_{12}X_2)
  +(-S_{11}^{-1}S_{12}X_2)^{T}S_{12}X_2
  +X_2^{T}S_{12}^{T}(-S_{11}^{-1}S_{12}X_2)
  +X_2^{T}S_{22}X_2
\end{align*}
\begin{align*}
  0<{}&X_2^{T}S_{12}^{T}S_{11}^{-1}S_{11}S_{11}^{-1}S_{12}X_2
  -X_2^{T}S_{12}^{T}S_{11}^{-1}S_{12}X_2
  -X_2^{T}S_{12}^{T}S_{11}^{-1}S_{12}X_2
  +X_2^{T}S_{22}X_2
\end{align*}
\[
  0<X_2^{T}S_{12}^{T}S_{11}^{-1}S_{12}X_2
  -X_2^{T}S_{12}^{T}S_{11}^{-1}S_{12}X_2
  +X_2^{T}S_{22}X_2
  -X_2^{T}S_{12}^{T}S_{11}^{-1}S_{12}X_2
\]
\[
  0<X_2^{T}S_{22}X_2
  -X_2^{T}S_{12}^{T}S_{11}^{-1}S_{12}X_2
\]
\[
  0<X_2^{T}
  [S_{22}-S_{12}^{T}S_{11}^{-1}S_{12}]X_2
\]
Sigue que
\[
  S_{22}-S_{12}^{T}S_{11}^{-1}S_{12}>0
\]
demostrando la suficiencia.

Suponga que el complemento Schur es cierto, se definen las matrices
\[
  A=[S_{11}-S_{12}S_{22}^{-1}S_{12}^{T}]^{-1}
\]
\[
  C=[S_{22}-S_{12}^{T}S_{12}^{-1}S_{12}]^{-1}
\]
\[
  B=-S_{11}^{-1}S_{12}C
\]

Desarrollando $B$
\[
  B=-S_{11}^{-1}S_{12}C
  =A[A^{-1}S_{11}^{-1}S_{12}C]
\]
\begin{align*}
  =-A[{}&(S_{11}-S_{12}S_{22}^{-1}S_{12}^{T})
  S_{11}^{-1}S_{12}
  (S_{22}-S_{12}^{T}S_{11}^{-1}S_{12})^{-1}]
\end{align*}
\begin{align*}
  =-A[{}&(S_{11}S_{11}^{-1}S_{12}
  -S_{12}S_{22}^{-1}S_{12}^{T}S_{11}^{-1}S_{12})
  (I-S_{22}^{-1}S_{12}^{T}S_{11}^{-1}S_{12})^{-1}
  S_{22}^{-1}]
\end{align*}
\begin{align*}
  =-A[{}&S_{12}(I-S_{22}^{-1}S_{12}^{T}S_{11}^{-1}S_{12})
  (I-S_{22}^{-1}S_{12}^{T}S_{11}^{-1}S_{12})^{-1}
  S_{22}^{-1}]
\end{align*}
\[
  =-AS_{12}S_{22}^{-1}
\]

Sigue que
\[
  S
  \begin{bmatrix}
    A & B\\
    B^{T} & C
  \end{bmatrix}
  =I
\]
Por lo tanto
\[
  S^{-1}=
  \begin{bmatrix}
    A & B\\
    B^{T} & C
  \end{bmatrix}
\]
y $S$ es no singular.